\section{Graphs, from scratch}\label{app:graphs}

The few terms of graph theory that the note uses, each with the place where it is used. The terms
of the study of the groups, shape and order of a group and the configuration model, are in
the note \emph{The structure of groups in free flight}.

\begin{itemize}
  \item \textbf{Graph.} A set of \emph{nodes} and a set of \emph{edges}, each edge joining two
        nodes. The active graph $G(t)$ of Section~\ref{sec:configurations} has a node for each flight
        and an edge for each pair interval in progress.
  \item \textbf{Path.} A sequence of nodes, each joined by an edge to the next. Two nodes are
        \emph{connected} when a path joins them.
  \item \textbf{Connected component.} A largest set of nodes connected to each other. Every node is
        in exactly one component; a node with no edge is a component alone (a pilot flying solo).
        A configuration is a component with at least two nodes, with its edges.
\end{itemize}

Two members of a configuration need not share an edge: in the chain 1--2--3 of
Figure~\ref{fig:graph}, pilots 1 and 3 are in the same configuration through pilot 2.

\begin{figure}
\centering
\begin{tikzpicture}[x=1cm,y=1cm, node/.style={circle, draw=ink2, fill=white, minimum size=6mm, inner sep=0pt,
    font=\footnotesize}, gedge/.style={ink2, line width=0.9pt}]
  \fill[blueI!10] (-0.5,-0.5) rectangle (2.5,1.5);
  \fill[blueI!10] (3.2,-0.5) rectangle (4.8,1.5);
  \node[node] (a) at (0,0) {1}; \node[node] (b) at (1,1) {2}; \node[node] (c) at (2,0) {3};
  \node[node] (d) at (3.5,0) {4}; \node[node] (e) at (4.5,1) {5}; \node[node] (f) at (6,0.5) {6};
  \draw[gedge] (a) -- (b) -- (c); \draw[gedge] (d) -- (e);
  \node[note] at (1,-0.85) {configuration $\{1,2,3\}$ (a chain)};
  \node[note] at (4,-0.85) {$\{4,5\}$};
  \node[note] at (6,-0.85) {solo};
\end{tikzpicture}
\caption{The graph at one instant: three components, two of them configurations.}
\label{fig:graph}
\end{figure}

\begin{itemize}
  \item \textbf{Directed graph.} The edges carry an arrow. For the episodes the nodes are
        groups, with an arrow $g\to g'$ when $g'$ starts as $g$ ends and they share a member
        (Section~\ref{sec:episodes}). An episode is a connected component of this graph, arrows
        ignored.
  \item \textbf{Bipartite graph.} Nodes of two kinds, edges only between kinds. At one instant, the
        groups ending and those starting, with an edge when they share a member. Each
        component is one event (Figure~\ref{fig:events}); its outcome follows from how many
        groups it has on each side: one to one is a join or a leave (members added or
        removed), one to several a split, several to one a merge (Section~\ref{gs-sec:describe} of the note on the structure of groups).
\end{itemize}

\begin{figure}
\centering
\begin{tikzpicture}[x=1cm,y=1cm, conf/.style={rounded corners=2pt, fill=blueI, text=white, font=\footnotesize,
    minimum height=5mm, inner xsep=4pt}, link/.style={->, ink2, line width=0.8pt}]
  \node[conf] (c1) at (0,0.6) {$\{1,2\}$}; \node[conf] (c2) at (0,-0.6) {$\{3,4\}$};
  \node[conf] (c3) at (2.6,0) {$\{1,2,3,4\}$};
  \draw[link] (c1.east) -- (c3.west); \draw[link] (c2.east) -- (c3.west);
  \node[note] at (1.3,-1.3) {merge: two to one};
  \node[conf] (d1) at (6,0) {$\{1,2,3\}$};
  \node[conf] (d2) at (8.4,0) {$\{1,2\}$};
  \draw[link] (d1.east) -- (d2.west);
  \node[note] at (7.2,-1.3) {one to one, 3 removed: leave};
  \node[note] at (0,1.3) {ending}; \node[note] at (2.6,1.3) {starting};
\end{tikzpicture}
\caption{Two events: each is a connected component of the bipartite graph of the groups
ending and starting at one instant.}
\label{fig:events}
\end{figure}
